% !TeX root = ../教材题库_全部条目_转写工作稿.tex
% Source: XJ-B1-C05-EX-P157-Y1411
% 第5章 / 5.1 任意角和弧度制 / 教材 p.157 / PDF 第164页
% OCR未校对：公式与文字均可能有误；未完成数学TeX转写；禁止直接发给学生
% 原题及相邻原解截图（核对用）：../题目原图/XJ-B1-C05-EX-P157-Y1411-01.webp
% 本文件只在转写完成后作为题干使用；OCR 条目尚不是可靠题干。
\begin{bankitem}{5.1 例题 例题2}{第5章 / 5.1 任意角和弧度制 / 教材 p.157 / PDF 第164页}{XJ-B1-C05-EX-P157-Y1411}{OCR未校对：公式与文字均可能有误；未完成数学TeX转写；禁止直接发给学生}
\begin{ocrdraft}
[教材p.157 / PDF第164页]
\par
00写出终边在3，轴非负半轴上的角的集合．
\par
解在0。、360。内，终边在3，轴非负半轴上的角是90。，
\par
如何用集合表示出终
\par
则所有这些角的集合为
\par
边在3，轴上的角？
\par
\{到一90。十．360。，k\ensuremath{\text{[OCR:EUR]}}Z\}.
\end{ocrdraft}
\end{bankitem}
